\documentclass[11pt,a4paper]{article}
\usepackage[utf8]{inputenc}
\usepackage[T1]{fontenc}
\usepackage[spanish]{babel}
\usepackage[margin=2cm]{geometry}
\usepackage{array,booktabs,longtable}
\usepackage{xcolor}
\usepackage{enumitem}
\usepackage{hyperref}
\usepackage{lmodern}

\hypersetup{colorlinks=true,urlcolor=blue,linkcolor=black}
\setlength{\parindent}{0pt}
\setlength{\parskip}{5pt}
\setlength{\emergencystretch}{3em}
\newcolumntype{P}[1]{>{\raggedright\arraybackslash}p{#1}}
\newcommand{\interaccion}[1]{\textbf{I#1}}

\begin{document}
\begin{center}
{\LARGE\textbf{AeroAndes Reservas: notas de presentación}\par}
\end{center}

\begin{center}
\small Repositorio del equipo: \url{https://github.com/iMouraad/APPD_MPZ}
\end{center}

\section*{Objetivo de la presentación}

Explicar cómo el diseño de AeroAndes combina comunicación síncrona y asíncrona para evitar la sobreventa de asientos, procesar pagos, enviar notificaciones y mantener el embarque durante una desconexión.

\subsection*{1. Vista general del sistema}

El pasajero puede usar la aplicación web o móvil y una agencia de viajes puede usar su propio sistema. Ambos canales llegan al \textbf{Servicio de reservas}. Este servicio coordina inventario, pagos y el estado de la reserva; no administra por sí solo la disponibilidad de los asientos.

La arquitectura se divide en cuatro zonas: canales de venta, centro de datos de AeroAndes, sistemas externos y red local del aeropuerto.

\end{document}
